% Summarize main question
This thesis thesis has explored the question of how to evaluate, train, and benchmark machine learning models for decision making problems. \

% Summarize BPR
In Chapter 3, we demonstrated how decision-aware model selection can lead to better public health outcomes.
We also introduced a new evaluation metric called Best Possible Reach which quantifies the decision quality of Top $K$ allocation problems.

% Summarize DAML
This work was extended in Chapter 4 where we demonstrate how a model can be \emph{trained} to maximize both decision quality and probabalistic maximum likelihood, in addition to answer the question of \emph{how to rank} locations based on predicted scores. \
The proposed Decision-Aware Maximum Likelihood method was compared to a small number of other decision aware methods, and found to be superior on the given tasks. \
We also explored problems beyond public health, including a wildlife conservation task.

% Summarize benchmarking
In Chapter 5 we turned our attention to understand why some decision-aware learning methods perform better than others, and to provide practical guidance for practitioners looking to implement these methods. \
Here we decomposed the analysis of decision-aware methods into two components: what solutions does the method prefer, and which solutions is the method capable of learning via gradient descent. \
We also found that many tasks used in common decision-aware benchmarks do not actually benefit from decision-aware methods, and propose several open-source alternatives. 

% Summarize thesis and suggest future benchmarking working
Taken together, these contributions provide end-to-end guidance for the training and evaluation of decision-aware machine learning techniques.
However, crucial open questions remain. Although we have demonstrated the importance of more careful benchmarking comparisons between methods, we still lack an analytical understanding of when decision-aware learning is beneficial, and how this depends on the structure of the downstream optimization problem and the role predictions play in determining the final decision. \
Future work should therefore develop new benchmark tasks that more clearly expose the conditions under which decision-aware learning offers advantages over decision-unaware two-stage approaches. \
Aditionally, new training methods are needed that make more principled use of prediction error within decision-aware learning as suggested in Chapter \ref{chap:daml}. \
Rather than discarding predictive accuracy as irrelevant, these methods should identify which errors are decision-critical and allocate learning effort accordingly, avoiding overemphasis on errors that are either too easy to correct or unlikely to affect the downstream decision. \

% suggest DTR, patient intervention, and inference
Beyond the largely social-science focused problems considered in this thesis, decision-aware methods have widespread applications in many domains. \
One particularly exciting domain is that of machine learning training and inference itself. \
In the twilight of Moore's Law, the rise of expensive foundation model inference has made computational efficiency more important than ever. \
Recent work has shown that by simply sampling more model outputs at test time, performance (as measured by at least one sample getting the correct response) can improve log-linearly with compute \citep{brownLargeLanguageMonkeys2024}.\
This has trasformed the simple act of using a model to output predictions into a budgeting problem. \
How we spend resources at test-time is a meaninful axis of model improvement: \cite{muennighoffS1SimpleTesttime2025a} showed consistent improvement over several-orders of magnitude, taking models from 0 to 60\% accuracy on challenging math datasets, and outperforming pretrined models with many more parameters. \
Here decision-aware methods can help answer the question of how to best spend fixed computational budgets to achieve optimal results from model inference. \
Furthermore, if model inference is now a decision task, should we be training all models in a budget conscious, decision aware way? \
We hope that the methods and findings of this thesis help to explore these questions in future work.